% ============================================================
%  USPSC-Cap3-Metodologia_3.1-Caracterizacao_da_Pesquisa.tex
%  Seção 3.1 — Caracterização da Pesquisa
% ============================================================

\section{Caracterização da Pesquisa} \label{sec:caracterizacao_pesquisa}

A presente pesquisa pode ser caracterizada segundo quatro dimensões complementares: sua natureza, seus objetivos, a abordagem ao problema e os procedimentos técnicos empregados.

Quanto à sua \textbf{natureza}, trata-se de uma pesquisa \textbf{aplicada}, uma vez que visa produzir conhecimento direcionado à resolução de um problema concreto — a ausência de protocolos metodológicos replicáveis para análise de sentimento no domínio financeiro em língua portuguesa — com potencial de aplicação prática por analistas, gestores e pesquisadores da área \cite{galdi2018}.

Quanto aos \textbf{objetivos}, a pesquisa é \textbf{descritiva e exploratória}. Descritiva porque busca caracterizar a distribuição de polaridade e os padrões semânticos presentes no corpus de publicações coletadas, identificando regularidades temporais e temáticas sem estabelecer relações de causalidade. Exploratória porque investiga um domínio ainda
pouco estudado na literatura nacional — a análise de sentimento de publicações financeiras em português brasileiro via modelos de linguagem pré-treinados — e porque a avaliação comparativa das abordagens de classificação gera hipóteses para pesquisas futuras \cite{santos2023finbert}.

Quanto à \textbf{abordagem ao problema}, trata-se de uma pesquisa \textbf{quantitativa}. O corpus textual coletado é tratado computacionalmente para gerar representações numéricas de polaridade (positiva, negativa ou neutra) que, agregadas, produzem um índice de sentimento passível de análise estatística descritiva. Embora a interpretação dos resultados recorra ao contexto qualitativo de eventos macroeconômicos do período, a unidade central de análise é de natureza quantitativa \cite{kearney2014}.

Quanto aos \textbf{procedimentos técnicos}, a pesquisa adota o modelo \textbf{experimental-computacional}, caracterizado pela construção, aplicação e avaliação de um \textit{pipeline} de PLN sobre dados coletados programaticamente via interface de programação de aplicações (API). Esse modelo é consistente com a tradição metodológica da área de PLN aplicado a finanças, na qual a replicabilidade e a documentação pormenorizada dos procedimentos são requisitos centrais para a validade científica do trabalho
\citeonline{santos2023finbert}.

O \autoref{tab:resumo_pesquisa} sintetiza as escolhas metodológicas adotadas e suas respectivas justificativas.

\begin{table}[htb]
	\caption{Síntese da caracterização metodológica da pesquisa}
	\label{tab:resumo_pesquisa}
	\begin{tabular}{p{3.5cm}p{3.5cm}p{8cm}}
		\hline
		\textbf{Dimensão} & \textbf{Classificação} & \textbf{Justificativa} \\
		\hline
		Natureza & Aplicada & Geração de protocolo metodológico replicável para o domínio financeiro em PT-BR \\
		\hline
		Objetivos & Descritiva e exploratória & Caracterização de padrões de sentimento; investigação de domínio pouco estudado na literatura nacional \\
		\hline
		Abordagem ao problema & Quantitativa & Mensuração numérica de polaridade via PLN e construção de índice de sentimento \\
		\hline
		Procedimentos técnicos & Experimental-computacional & Construção, aplicação e avaliação de \textit{pipeline} de PLN sobre dados coletados via API \\
		\hline
	\end{tabular}
	\legend{Fonte: Elaborado pelo autor.}
\end{table}
